\documentclass[a4paper]{article}
\usepackage{graphicx}
\usepackage{amsmath,amssymb}
\usepackage{hyperref}
\usepackage{minted}
\usepackage{multirow}
\usepackage{float}
\usepackage{todonotes}
\usepackage{forest}

\title{}
\date{}

\begin{document}
    \maketitle
    The simplest Transformer performs sequence modeling through a mechanism called \textbf{attention}. Instead of processing elements strictly one after another, it allows every element in a sequence to directly consider all other elements.

    Assume a sequence
    
    \[
    x_1, x_2, x_3, x_4
    \]
    
    Each \(x_i\) is a numerical vector (for example, sensor measurements at a specific time).
    
    For every element \(x_i\), the Transformer computes three vectors:
    
    \[
    q_i = W_q x_i
    \]
    
    \[
    k_i = W_k x_i
    \]
    
    \[
    v_i = W_v x_i
    \]
    
    These are called
    
    \[
    \text{Query (Q)}, \quad \text{Key (K)}, \quad \text{Value (V)}
    \]
    
    where \(W_q\), \(W_k\), and \(W_v\) are learned weight matrices.
    
    For example, for \(x_3\):
    
    \[
    q_3 = W_q x_3
    \]
    
    \[
    k_3 = W_k x_3
    \]
    
    \[
    v_3 = W_v x_3
    \]
    
    The model then determines how much \(x_3\) should attend to every element in the sequence. This is done by comparing \(q_3\) with all keys:
    
    \[
    q_3 \cdot k_1,\quad q_3 \cdot k_2,\quad q_3 \cdot k_3,\quad q_3 \cdot k_4
    \]
    
    These values are called \textbf{attention scores}.
    
    The scores are then normalized using the softmax function to produce weights. For example:
    
    \[
    [0.1,\; 0.6,\; 0.2,\; 0.1]
    \]
    
    This means that the third element \(x_3\) pays the most attention to \(x_2\).
    
    The output representation of \(x_3\) becomes a weighted combination of the value vectors:
    
    \[
    0.1 v_1 + 0.6 v_2 + 0.2 v_3 + 0.1 v_4
    \]
    
    Therefore, each position in the sequence aggregates information from all positions, weighted according to their relevance.
    
    In matrix form, the attention mechanism is written as
    
    \[
    \text{Attention}(Q,K,V) =
    \text{softmax}\!\left(\frac{QK^T}{\sqrt{d}}\right)V
    \]
    
    where \(d\) is the dimensionality of the key vectors.
    
    Conceptually, the difference from an LSTM is the following.
    
    An LSTM processes a sequence sequentially:
    
    \[
    x_1 \rightarrow x_2 \rightarrow x_3 \rightarrow x_4
    \]
    
    Information must propagate through each intermediate step.
    
    A Transformer instead allows direct interactions between all elements:
    
    \[
    x_1 \leftrightarrow x_2 \leftrightarrow x_3 \leftrightarrow x_4
    \]
    
    The attention mechanism determines how strong each connection should be.
    
    For time series data this is particularly useful. If a pattern at time \(t=50\) resembles something that occurred at time \(t=3\), the Transformer can directly assign attention to that earlier position. In contrast, an LSTM must carry that relationship through many intermediate recurrent steps.

    % \bibliographystyle{plain}
    % \bibliography{/home/donkarlo/Dropbox/repo/shared_research/bibliography/bibliography.bib}
\end{document}